% 第16章 D) 被埋葬者们;原书p397-412

\chapter{D)~被埋葬者们}

\section{X.~灵车}

\subsection{棺材1——或知恩的D-模}

\notehead{93}{}

\dmark{}(5月21日)我的反思在那些“根正苗红”的学生们——“1970以前”的——身上,已滞留两个星期。每一天,这场反思都摆出一副“最后一项补充”的面目,出于求个心安,补向一场(几乎)已告终结的反思。不止一次,是一条不起眼的脚注——贸然接上前一天或前两天的反思——越拖越长,直到一份独立“注”的规模。每一回,它都很快得名,把自己区别于所有其他的注,并恰好插进其葬礼行列的合适位置,仿佛它从来就在那里!每隔两天,我都得(每次都怀着乐趣)至少重做目录的结尾:目录仿佛已收尾,却随之添上两三位新的入列者,有时甚至一整个新的行列……

这支队伍的规模,渐渐变得令人不安:谁也不会肯把这全部读完!但它这样加长,老实说,并不是为了一个假想读者的可疑收益,而首先是为我自己的收益——一如我做数学之时。这些“最后一项补充”——我每一回都仿佛违背本意地登船——从不曾让我后悔。借着一次次的最后一项补充,我学到了许多别处无从学到的东西,倘若不这样“从材料出发”反思,我便无从学到这些。而这些东西,一件一件聚成一幅色彩鲜明的图景,宏阔而有多个侧面。此刻我仍看到它尚未全部完成——还有两处,似乎在讨最后一笔。

我觉得,谈过那些“根正苗红”的学生之后,如今也该稍稍谈谈\textbf{被埋葬者们}了——那些“与我一同,有权领受这场以沉默与轻蔑进行的下葬之殊荣”的人。一如我、一如那些兴冲冲下葬的人,这些被埋葬者也不是圣人,也没有殉道的使命。我相信,他们中没有一个不曾暗暗怪我\dmark{}\origpage{396}给他招来麻烦——那些我全然无意给他招来的麻烦(只因为他曾失于轻率,把赌注押在我身上,押在某种做数学的进路与某种做派上……)——或者,至少不曾试图与我划清界限,一旦认清赌注分明已输\fnote{1}{(*)}{(1985年2月)自我到蒙彼利埃以来,在校外,我总共得知七八份(简短的)出版物,呈献(以摘要的形态)一项与我一起做、受我启发的工作。我的名字,在所有这些里缺席。}\origpage{397}。何况我可以确认:那是白费力气——一旦被盯上,就完了;划清界限,只会喂养一份轻蔑、给它一份心照不宣的辩解,而不能解除它。同样不止一次、以多种方式,我看到下葬者与被葬者的角色相互毗邻、相互混淆\fnote{2}{(**)}{(9月2日)以各不相同的方式,他们每一位都曾在某个时刻,终于把对其工作的轻蔑内化并据为己有,默许那份藏掖其工作或把它归为“无趣”的共识。}。想必正是这些暧昧的侧面,造成我身上一份长久的迟疑:不肯比此处彼处顺带的影射更周详地谈“被埋葬者们”。很可能,除Zoghman之外,另外三位我认识的人里,没有一个会感激我在这里替他做“广告”——仿佛我给他招的麻烦还不够多似的。

一如《收获与播种》里的许多时刻,我终于越过我身上这份迟疑。我对自己说:即便对着那些曾因我而受苦的人——经由他们在某个时刻做出的选择;在那选择里,他们出于这样或那样的缘由,分明各得其所,尽管他们同我一样,不曾料到其选择所附带的种种不利——即便对着他们,我的角色,也不是帮他们回避一个实实在在的局面;无论愿意与否,他们都身陷其中,而这局面分明有其意味,哪怕它带着严重的麻烦。

在接上我那些亲爱的共同亡故者、共同被埋葬者的四口棺材的黑色系列之前,我或许该以一条不那么丧气的注,给读者换换心情。首先,在与我们大学数学研究所的“地方”层面的关系里,我完全没有体验到:我为一个职位的候选人所说的话、或候选人是我的学生(不言而喻,是“1970以后”的)、或其著作受我著作的影响,就必然于他不利。那种系统性抵制的态度,仅仅是“上流”数学界对我这个人的关系——推而广之,对那些显得与我“1970以后”相连的人——的特征。据我所知,这份抵制,自我离去以来的十四年里几乎毫无破绽,只除两桩寒微的例外。其一涉及一个学生:他开局颇有希望,本该跟我做一篇题材极为诱人的国家博士论文;他申请我们大学的助理教授职位,却被专家委员会驳回。他在全国一级被“捞了回来”,有Demazure的帮助——我曾就这位学生的工作给Demazure写过信\fnote{3}{(*)}{就晋升或取得职位与身份的“实际”层面而言,我1970年以来的教学活动的全部账目,归结为两桩取得职位并附带身份的事:一回是助理教授,另一回是助教。说来奇怪的讽刺:两回,这晋升都是当事人一切研究活动戛然而止、彻底停摆的信号。}。另一方面,有两次,《Topology》杂志接受了出自我的学生的文章:一篇“Stein因子分解与剖分”,出自Jean~Malgoire与Christine~Voisin;以及Yves~Ladegaillerie一篇将刊的文章,包含其1976年论文的中心结果(见注n\textsuperscript{o}~94)。

我尤其已有机会谈到过Zoghman~Mebkhout,这里只“存照”再谈\fnote{4}{(**)}{除《导言(6)》(葬礼)之外,谈到Mebkhout的还有:“我的孤儿们”“当值无名之辈与好上帝定理”“不义——或归来的意义”“反常性”“墓外相会”“受害者——或两种沉默”“大部头与上流社会”“赊账的学位论文与全险”各注(n\textsuperscript{o}~46、48$'$、75、76、78、78$'$、80、81)。}。Mebkhout想必自1974年起开始从我的著作汲取灵感,并且顶风冒雨、一直汲取到今天。我不知道我的哪一位“官方”学生做出了分量可比的著作——尽管Mebkhout的著作,必然带着其被迫在逆境中推进的印记。如我在《导言(6)》里所说,四年来,Mebkhout的观念与结果,人人都在用,而他的名字却被小心地藏掖\fnote{5}{(***)}{在那场下葬里充当掘墓人的,大有人在;1981年6月的吕米尼研讨会,几乎全员参与。除我的上同调学家学生们(对此可参见“我的学生们(2):团结”一注,n\textsuperscript{o}~85)之外,其职业诚意在此直接且严重地成为问题、而我有知情的,是 J.L.~Verdier、B.~Teissier、P.~Deligne、A.A.~Beilinson、J.~Bernstein。}。我的朋友何以能在承受轻蔑、继而承受不义——如一场不可抗拒的宿命——的同时继续做数学,这对我是谜——那份宿命,经由他想必(并且想必仍然)感到高悬其上、令人眩晕的人而来\fnote{6}{(*)}{当然,Zoghman~Mebkhout并不比我傻,而且他足够“在局中”,对我每一位上同调学家学生的工作都有精确的把握,能掂量其分量与限度,毫无理想化的倾向。但这并不妨碍:一些力量可观的抑制,使他连那样一个念头都不曾有——即公开质疑他们中的任何一个,哪怕在恶意分明之处。}\dmark{}\origpage{398}:那些人,他想必头一回听说时,俨然“赛场诸神”;而那时,他(一如当年的我)是一个生计拮据的寒微留学生。1979年答辩之时,他在奥尔良有一个助教职位。当时他竭力设法进CNRS,三次进攻——到第三回(1982年10月),人们终于肯给他一个研究员职位(相当于大学的助理教授或副教授)。这给了他,即便不是身份的担保,至少一份相对的安全。

在我所知的四位“共同被埋葬”的数学家之中,Mebkhout是唯一一个逆着所有人继续其工作的:他信任自己的数学本能,不被那份无情时尚所能提示的谨慎与机会主义的考虑绊住。他身上——他并非好斗的天性——有一种对自己判断的素朴的\textbf{信念},它同时也是一份\textbf{宽宏};而它(远甚于脑力“手段”)是做出革新而深刻的工作的首要条件。我对他的著作的看法,肯定仍不完全。就我所知其著作的主干部分而言,我觉得:以他那出众的手段,倘若处身于一份热忱而起作用的同情的气候里,他本可以在三年或四年里(而非十年)完成它,并把它带向更大的成熟——在喜悦而非辛酸之中。但三年也罢十年也罢,“成熟”与否,值得注意的,是革新的著作出现了,而且竟能在这样的条件下出现。

\subsection{棺材2——或电锯剖切}

\notehead{94}{}

Yves~Ladegaillerie1974年开始与我工作。那是“姑且一试”,在他生涯的一个低谷——我当时向他提出了几条关于1维拓扑复形在曲面中的嵌入的天真反思;那时我对曲面一无所知(除亏格观念),他懂得更少。这有点像Grothendieck学问(反正在我这里,一切总是这样开始……);它在他那里或多或少咬住了钩,直到有一天终于“翻了盘”——我已记不清何时何故。也许是在一个分明多汁的问题正在成形之时:关于紧致带边定向曲面中1维紧复形的同痕类的判定的一条关键猜想。真——假?悬念持续了六个月、一年;其间伊夫熟悉了(也顺带使我熟悉了)曲面理论的关键定理,同时推进其工作的“奠基”部分。已知的结果使猜想相当可信,却分明远远不够——因为那猜想蕴含巴厄与Epstein的那些难啃的结果,以及另一些有着离奇甚至可疑面目的东西。1975年夏,他终于证明了那条关键猜想。它实质上等价于:以基本群的语言,对一个可三角剖分的紧空间嵌入紧致带边定向曲面的全部同痕类,给出完整的代数描述\fnote{7}{(*)}{非定向情形的“类似”陈述是假的——分明这是一个精微的结果,被小心地“剖”在一组同样“可信”却仍为假的假设-结论之中。关于Ladegaillerie工作的其他评语,见《纲领概要》,尤其第3节开头。}。

\dmark{}从伊夫“咬钩”那一刻起,他一年、至多一年半做完了论文:结果、成文,一应俱全,而且几乎苛刻地完美。那是一篇光辉的论文,比我多数学生的论文薄,却与其余十一篇中的任何一篇同样实。答辩在1976年5月举行。

论文至今未出版。它虽不厚,据说还是太厚,不宜出版——这是人们给我的许多条极好的理由之一。我在“进步是挡不住的!”一注(n\textsuperscript{o}~50)里列举了几条。我为“安放”这篇倒霉论文所作努力的经过——它是我有幸激发的最好论文之一——可以写一本小书;它肯定富有教益,而我放弃写它。在那些有极好的理由忘掉去了解这些结果、并闭着眼睛把这一切埋葬的昔日好友之中,有(按登场的顺序)Norbert~A.~Campo、Barry~Mazur、Valentin~Poenaru、Pierre~Deligne——还没算经Springer公司转手的 B.~Eckmann\fnote{8}{(*)}{我与Eckmann并不相识;为使伊夫的论文在Lecture~Notes出版而做的通信,是与Springer主管LN的Peters博士进行的。我想,经由六十年代由我(SGA尤其是)或由学生(论文)出版的十五卷LN,我曾是凭我的担保为这套丛书尚在草创之际的前所未有的信誉与成功出过力的人之一。拒绝我所推荐的工作所给的理由(他们不出版论文),是句玩笑话。通信事务上的“新面貌”,我的头一份经验也出自这一情节:怀着一套真正令人叹为观止的默契, A.~Campo、B.~Mazur、 V.~Poenaru与Peters博士,对我第二封信,都竟免于赏我一复——当时我天真(我的悟性慢……)地卷土重来,而他们那封迟疑的回信已表明:他们没有费心了解Ladegaillerie工作导言里所陈述的结果。}\dmark{}\origpage{400}。中心结果终于要发表了——九或十年之后,削减到只剩骨头,以《Topology》一篇短文的形式(嘘——我在那份可敬杂志的编委会里有一位同谋……)。工作的其余部分,一边证明了人人久已不证而用的事情(当然,没有证明也过得很好!),一边发展典型的Grothendieck学问,全然背逆时下的惯例与良俗。我很清楚:倘若我的朋友Deligne不在未来十年里大张旗鼓地负责“发现”它们,别人在三五十年内也忍不住会重做它们——因为我健全的本能告诉我:那些是基本的东西。它们曾是我阿那贝尔式沉思中宝贵的引线;倘若上天借我余年,在《数学反思》发展阿那贝尔式代数几何瑜伽的那一部分里,我会有充分的机会指涉它们。

这段经历对我是一次显影,同类中的头一次——对某样东西的显影;那东西,直到葬礼的反思,我才充分认识。况且,自从我的心神另有所系,我倾向于把它忘了。Yves~Ladegaillerie——我有过的最光彩的学生之一——则从那一刻起就懂得了:要在今天的数学界被接纳,仅仅全力投入、并做出合乎一切卓越要求的工作,是不够的。他另有几手,七年来他投身于更接地气的任务与较少问题的回报。他运气好:早在与我那场不幸相遇之前,他就有助理教授的终身职位\dmark{}\origpage{401},给了他一份其厄运不曾危及的保障。去年,一点数学的火星似乎重新燃起,在一个与我近年来关心的问题十分贴近的主题上——瑟斯顿式的双曲几何及其与Teichmüller群的关系。甚至可能我们还会一起走一小段路;或者他走他个人的散步,纯粹为了乐趣,不期待数学本身所能给的回报之外的任何回报。他清楚:倘若他期待别的,他最好换一个对话者或同路人(甚至,换一段过去……)。

\subsection{棺材3——或有点太相对的Jacobi簇}

\notehead{95}{}

我与Carlos~Contou-Carrère的最初相遇,是在1973年我到蒙彼利埃之后不几天,数学所的走廊里。他会把我堵在某个昏暗的角落,在我来得及客气告退之前,把一股数学解释的洪流倾倒在我头上。他连珠炮般倾倒给我的东西,全从我头顶掠过;而他丝毫不曾察觉,或在我怯怯暗示之时有丝毫窘迫。他迫切需要对话者,而我并非他唯一的“违心对话者”。那又是一个我全然不沾数学的时期。有一两年,我一看见他(一眼就能认出)的身影在走廊尽头出现,就躲。直到林登——他曾在蒙彼利埃做过一年客座教授——向我透风:Contou-Carrère手段不凡,却眼看要触礁,因为不知道怎样使用它们。在那之前,Contou-Carrère倾倒给我的东西站不站得住、他有没有手段,连掠过我心头都不曾——一切那么遥远。也许林登的提示来得恰逢其时:我开始对数学问题重新发生一点兴趣。总之,我横下心来,问Contou-Carrère可肯给我解释一件他做过的事,好让我能懂。我疑心:我是头一个向他提这种要求的人——至少在他来法国已有好几年之后。要他显豁一件事不容易,却绝非不可能,而且值得。我很快发觉:林登没有看错——Contou-Carrère满脑子观念,只等着被提炼、被细心发展;而且在我们可能向他提交的几乎一切数学局面里,他都有一种即时的、非常可靠的直觉。凭这份直觉的敏捷与可靠——哪怕在他毫不熟悉的事情上——他超过我、令我惊叹;在别的学生身上,我以可比的程度体会到它的,只有Deligne\fnote{9}{(*)}{我不确信在其他数学家身上遇到过它——除Pierre~Cartier(他年轻时以这份卓越的能力深深打动过我)与Olivier~Leroy(下一条注将谈到他)。}。相反,他对书写有一道几乎完全的阻塞!不可思议:他做数学\textbf{不写}——天知道他怎么能哪怕做出那么一点,更不必说与他人的交流:在那方面,“触礁”是彻底的(见上文)。

倘若我有什么紧急而有用的事情要教给Contou-Carrère,那是书写的艺术——或更直截了当地,哪怕只是让他明白:数学是\textbf{写着}做的。我费了两年、也许三年,直到1976或1977年\fnote{10}{(**)}{(6月7日)核实之后:直到1978年2月。},不敢说是否全然成功。他第一部完全白纸黑字写成的大工作,是他的Schubert闭链论文——迟至去年12月(1983年)才答辩\fnote{11}{(***)}{那是一项长工作(我没有读过):他细心发展一些与我毫无关系的想法,其中给出一切“Schubert”型闭链的显式奇点消解——此前无人能做到。难得他有一次成形的写作,人们却责备它太详(更不必说他的陈述太一般……)!至于我,倘若我有批评,那将朝着相反的方向:虽然Contou-Carrère声称其方法适用于一切类型的半单群与Schubert闭链,他却只在一般线性群的情形做了工作——因此他在这件明确的事情上没有做到底:普适Schubert闭链的等变奇点消解\textbf{以及}这些Schubert闭链的奇点轨迹的描述。这份欠缺,在我看来,是那道对“从材料出发”的扎实工作、对书写的“阻塞”的一份遗产;那阻塞,长久以来是他的主要缺陷。}。1978年至今,我们的关系极为零星;我的角色,实际仅限于在他职业生涯中屡屡以这样或那样方式陷入困境之时——那职业,始终悬在最不牢靠的代理助教职位上——尽力扶持他。

有两三年,我曾试图向Contou-Carrère提供一套精确而柔顺的数学语言的基础,外加几条系统化的原则。带着这份行装,连同他的手段与观念的丰富,接续什么,他有真正的挑选余地。他没有从他自己的出发观念入手,而是接上了相对局部与整体Jacobi簇的理论——我曾把它作为可能的论文题目向他提过。一旦我放手让他自己干,他在仅仅一年之间做出了一项很漂亮的工作,其中一部分在一篇科学院报告短文里已有预告(注95\textsuperscript{1})。把这个矿脉做到底,本会代表几年令人着迷、强烈激励他的工作——足以同时学得概形技术的全部精微。那时我毫无疑心——明摆着,Cartier、Deligne、Raynaud三人都会热烈接待这项已完成的工作:深刻、艰难,且在好几个方面出人意料。Cartier确实满心欢喜地看到他的某些旧观念取得新的现实性。相反,是Raynaud的冷漠,一如Deligne的——后者把完整手稿在抽屉里压了六个月,不屑给一点活着的音讯\fnote{12}{(*)}{Contou-Carrère倒也先防了一手:在其短文里,对我的这个人——曾给了他出发纲领的人——只字不提。那也是白费:任他往里加多少自己的东西,有一种“做派”不会错,它系于某些主题;谁愿在今天的数学里混出前程,最好避开这些主题。(6月7日)向当事人查询之后,我发现:我在这里把围绕Contou-CarrèreJacobi工作的两个不同情节弄混了。细节与精确出处,见下一条注(n\textsuperscript{o}~95\textsuperscript{1})。}。这是二对一——足以感到风向。那有点太相对的Jacobi簇,被无限期弃置,任其自生自灭。电锯干得很出色……

这并没有为Contou-Carrère挡住厄运;为其写一份周详的报告,本可以再写一本小书,我由衷放弃。我想正是在这时——自我1970年离开那个机构以来(在那机构尚无自己的屋顶的年月里,我曾是四年间1958--62唯一在“实地”代表它、使它取信于人的人\dmark{}\origpage{404}),唯一的一次:我担当起为某人推荐一份邀请(为期一年)的责任,正当Contou-Carrère眼看要没有职位、流落街头。我知道:我推荐的人——一如当年的Hironaka、Artin或Deligne被我热情接待时那样无名——会像他们一样,为接待他的机构增光。当然,我不曾忘了这么说。对Contou-Carrère幸运的是:他的代理助教职位(诚然配不上一份如此挑剔机构的邀请之荣)最终得以续聘\fnote{13}{(*)}{我无从抱怨:五六年之后,借去年IHES建所二十五周年庆典之机,人们倒真赏了我这个人一份邀请,甚至让我在部长致辞的庄严招待会与随后一周的IHES居留之间选择——还包食宿(人家郑重向我担保)。我对老朋友Nico~Kuiper说:想到我,非常客气;但我这把年纪,不再出门旅行了……}。

对这一情节我并不太惊讶:当时我已了解Deligne的心境;况且Nico~Kuiper曾提醒我,在这个案子里一切取决于Deligne。(我甚至没有想到向他暗示:这事兴许也关涉科学委员会的其他成员——恰恰因为这是个案……)然而,在Contou-Carrère(我的“门生”——Verdier曾在一封信里想当然地这样称呼他……)的一切厄运之中,最触动我的情节,发生在1981年10月:事涉他申请佩皮尼昂的一个教授职位。佩皮尼昂的同行们(他在那里任代理助教)想必乐于身边有一个在任何数学分支里都从容、可咨询的人。一次教授职位出缺,他们把他列为唯一候选人——罕见之至的事,分明表明:他们想看的正是他、而非别人。Contou-Carrère在其阿根廷Santalo门下的博士论文之外,出版物相对少:主要是科学院报告短文,预告结果(其中一些深刻),而无证明。还没有人向他暗示过:在如今的时势下,在安身之前,最好有附完整证明的文章作“证据”——这一点我这边当然向他念叨过不少,只是从较不功利的角度\fnote{14}{(**)}{前一年,Contou-Carrère申请雷恩的一个教授职位;他认识Berthelot与Larry~Breen。当时大学咨询委员会认为他的申请可受理;但职位给了另一位候选人。没有人费心提醒当事人:倘若他想有职位的希望,就得发表他所预告结果的详细证明。次年委员会的驳回,对Contou-Carrère、对他的佩皮尼昂同行们、对我,都来得完全意外。带着事后的认识、在这场反思的光照下,我进而怀疑:随着其论文的成文(如今已被宣告“照原样不可出版”)与答辩,局面是否真的变了;他是否有希望在法国找到一个教授职位。}。总之,Contou-Carrère的申请\dmark{}\origpage{405}被大学咨询委员会判定不予受理,材料退回。当时使我倒吸一口凉气的,是:委员会主席(做出决定的全国机构)不以委员会名义、任何成员也不以个人名义,具有这点最起码的尊重:给当事人Contou-Carrère本人、或至少给佩皮尼昂数学所所长写几行,解释这次表决的意味——在毫无解释的情况下,这次表决只能被领受为对佩皮尼昂同行们选择的响亮否定,对其唯一候选人能否体面胜任其被提名之职的否定。委员会里有我的三位旧学生,其中两位亲自认识Contou-Carrère。他们当然知道他同他们一样是我的学生——何况材料里还包括我一份对候选人工作格外称道的报告。他们中没有一个、委员会其他成员中也没有一个,想到过这记不经任何手续的铡刀式表决所代表的羞辱,想到过对一位不比他们中任何人逊色的数学家的按部就班的鱼雷攻击。

正是这一事故,在我做数学家的一生中,头一回使我感到那种我已在反思中不止一次谈过的“气息”。四年前,经由“外来者”的情节,我已感到过它\fnote{15}{(*)}{见“我的告别——或外来者”一节,第24节。}。但那不是在我曾经所属的世界内部,吹在\textbf{他们自己人}身上——吹在一个毫无保留地认同于那个世界的人身上。我好几个星期、也许好几个月为此如病如恙。为摆脱当时攫住我的一份焦虑——我当时无心去认识它\fnote{16}{(*)}{我是在次年一场漫长的禅观期间,才意识到这份焦虑的;在那场禅观里,我发现了焦虑在我生活中的角色——它的在场(1976年之前是慢性的,1976年之后是偶发的),曾是我一生的“世上守得最好的秘密”。有一套效率极高的机制,藏掖焦虑的一切公认的征兆;那焦虑,对我自己、对我的亲人,都始终不为人知。},我四处奔走,东一封信西一封信,还写了一篇三十来页的文本《大脑与轻蔑》,带着黑色幽默的笔调;我终于放弃了发表它\fnote{17}{(**)}{劝我别发表它的,恰恰是那些我正准备为之开战的人——幸亏我有先见之明,在试图让它面世之前,把我的文本给他们看过。}。带着事后的认识,我看清:那本是就所遇之事的意义展开\textbf{禅观}的绝好时机。最滑稽的是:当时“妨碍”我觉察需要一场深入禅观的,恰是一场我那时正投入的长禅观\fnote{18}{(***)}{见“扫兴的老板——或高压锅”一节,第43节。}——而且是一场关于我与数学之关系的禅观(即便还不是关于我数学家过去的)!它被一个情节搅乱:生活以强力向我发问——而我在奔忙中回避那发问,随后重新扎进“禅观”。带着事后的认识,我看清:那场“禅观”当时配不上这个名字;它缺少真正禅观的一个本质维度:对\textbf{当下时刻}的我自身的留意。我当时“禅观”某些较为遥远的事件的意味,却忽视一份被压抑的焦虑(诚然,由于长期练习这类控制,它被完美地控制着)——那是我拒认那股被驳回的“气息”所带来的讯息的征兆。

但我正在离开本题。那场鱼雷攻击,当然有其必有的效果。佩皮尼昂的同行们被申饬了一次,就够了。显然,他们那里连代理助教的职位也没有了——至少对Contou-Carrère没有。他在蒙彼利埃临时找到一个替补职位,就本学年,编制者明年回来。

不过我并不太为他的前程担忧:一段时间以来,Contou-Carrère已有智慧抢在命运的打击之前,转而投向计算机。以他那出众的手段,想必久已高屋建瓴地驾驭这门学问,同时在闲余做他喜欢的数学。他是两个孩子的父亲;在如今的时势下,带着紧追不放的过去,数学分明是碰运气的事,甚至凶险。他大可以做一个辉煌的计算机学家:没人会因他多少当过我的学生而对他严苛。

\notehead{95\textsuperscript{1}}{}

(6月7日)大约77年底,我向Contou-Carrère提交了一份为相对局部与整体Jacobi簇理论的周详工作计划;其中包括,在局部情形,建议“重新拧紧”Jacobi簇与Cartier的ind-群,以找到一个具有更美泛性质的“完备”Jacobi簇,它应当是“自对偶”的。我没有任何可提出的证明,1978年2月之后也不再过问其工作:我发觉我的在场抑制而非刺激其才能。他在随后的一年里到底“启动”了;其第一篇注“相对曲线的广义Jacobi簇:构造与分解的泛性质”(论整体情形)发表于1979年7月16日(CRAS t.289, Série A - 203)。

次月,他为局部Jacobi簇找到了决定性的结果,却一年半不发表;然后在1981年3月2日的一篇科学院报告里发表“一半”(寻常的、未与Cartier群重新拧紧的相对局部Jacobi簇的泛性质),顶着(乍看不太有说服力的)名字“相对几何局部类域论”(CRAS t.292, Série I - 481)。至于完备局部Jacobi簇的理论——依我看更有趣味——存在一份科学院报告的底稿,从未发表,题为:“局部Jacobi簇,普适Witt双向量群与tame符号”。当然,1979年起我就知道他的结果:那是对我向他提议的临时纲领的一个完整实现;为它,须克服可观的技术困难,要求大量的想象力与技术力量。(如无差错)我只读过第一篇注;我奇怪他不发表后续,即局部部分——他从不曾就此清楚说明——但他分明对第一篇注所受的接待感到失望。1980年雷恩申请失败之后,鉴于附在他申请材料里的我的支持信\dmark{}\origpage{408}论及整体与局部相对Jacobi簇上的卓越结果,他想必还是认定谨慎为上(为准备次年佩皮尼昂的申请),至少再发表一篇关于局部Jacobi簇的注,如非和盘托出。又是两个月之后,1981年5月,他把第三篇注的底稿寄给Deligne与Raynaud(Cartier想必久已知情),我想是先探探路。(我不相信:自1979年8月他手握结果以来的任何时刻,他找Cartan呈送这第三篇注会有丝毫困难。)Raynaud与Deligne都不给他一点音讯——但1982年3月,Deligne给他寄来一篇文章的手稿“tame符号评注”,作者是加藤和也,献给Deligne;它在基底域的情形做出Contou-Carrère的理论,并猜想其对任意基底环有效。Contou-Carrère当时跟我谈起,自认深信Deligne把他的结果(不提其名,也不给证明的任何指示)通报了加藤。当时这事在我看来如此不可思议,以至于我没有把Contou-Carrère当真——如今我才明白:那完全是我光彩夺目的朋友Deligne惯常的“暂停!”做派。Contou-Carrère一副真正被冒犯的神气:有人竟“斗胆猜想”他俨然视为某种私产的东西。然而他自己的猜想得自我的传授,却同样不认为有必要在三条注的任何一条里影射我这个人\fnote{19}{(*)}{关于我在这类局面里、与我的某些学生之间常常扮演的那份共谋的角色,见“暧昧”一注,n\textsuperscript{o}~63$''$。}!他对我说来,想必理所当然;而Deligne对他做的单纯的动念却冒犯了他——但他不敢向当事人吐一个字。(我曾力劝他向Deligne说清;他却小心没有做……)

这些年来,我料想,他必定以某种方式对自己施加了暴力,不发表那些极美的结果——做它们时,他想必全力投入。倘若他这样对自己施暴,那是出于时势的考虑:分明不利于这一类Grothendieck学问。近几天他十分惊讶地收到同一位Deligne的一封信:他(不动声色地!)惊讶他没有发表关于“总体”Jacobi簇的那篇注,并向他索要他关于这一主题、甚至别的主题的一切。Zoghman~Mebkhout几天前已告诉我:Deligne正在使用这些东西,甚至在这一语境里点了Contou-Carrère的名。分明,时候快到了:Contou-Carrère终于将认下一个属于他的孩子——一个他谨慎地埋了近五年的孩子。也许,谁知道呢,两位“学生-仇敌”和解的钟点也已敲响:我的学生中最光彩的两位,一位是受勋的院士,一位是代理助教;然而(无论和解与否),久已是一对\textbf{兄弟}。

\subsection{棺材4——或没有鲜花也没有花冠的topos}

\notehead{96}{}

(5月22日)倘若我说我从未见过Olivier~Leroy,我几乎没有夸张。可以肯定的是:从他听说我那一刻起,他就决定像躲瘟疫一样躲我。他的缘由,我承认,向我滑脱。也许一个本能告诉他:我只会给他招麻烦;也许Contou-Carrère(他长久以来与他过从甚密)向他提过醒——我恐怕永远无从知道。但我到底有幸与他有过两次有实质内容的交谈;我记得很清楚。

第一次应当在1976、77年:我和Contou-Carrère不打招呼地登门看他,想聊聊数学——是否心里另有盘算,我说不好。不过,想必有言在先:奥利维耶考虑读第三阶段博士,而我袖子里有的是题目。因为在Contou-Carrère家见过他一两次,加上Contou-Carrère自己向我透露的,我有一个印象:奥利维耶悟性快,而且不只在数学上。三人那晚令人难忘。我很快向奥利维耶提起一个纲领:一个topos基本群的理论、以及topos框架里的Van~Kampen型定理;他一副感兴趣的样子。经由Contou-Carrère的代数几何研讨班,他想必有一点topos的常识;他似乎乐于有机会在一个具体理论的例子上“练练手”,用topos的语言。足足两三个钟头,我把我所见当发展的理论的周详总施工图倾倒给他;它随我谈论而渐渐丰满,同时,代数几何与拓扑的一大群具体局面在我身上浮起——那些局面,要在topos框架里表述;而每一回,我都得先向一个头一回听说的人“复述”\dmark{}\origpage{410}。那晚不止一次,Contou-Carrère(他几乎无所不读、胃口极好)眼神茫然、掉了队——连他也觉得一次太多;不止一次,我谨慎地问奥利维耶,是否今天到此为止、改日再续。我原可省了这份心——分明奥利维耶神清气爽,眼神明亮,完全自如;我甚至笑起来,因为难以置信他竟不垮——可他全然不垮!那是个也许二十岁的年轻人,只有一点概形的常识、一点拓扑与topos;我相信他倒也摆弄过不少无限离散群……老实说,那是三脚猫的底子;而凭这个,他竟能填上所有的空白,毫不费力地“感到”我这个老资格、凭十五年与这一主题的熟悉、在两三个钟头里全速讲述的东西。我从未见过这样的事——至多在Deligne身上,也许还有Cartier;Cartier年轻时在这一条线上也相当出众。

总之,分明是拍板成交:奥利维耶将就上述题目做他的第三阶段博士论文。然而他想必不曾料到路的尽头有什么。总之,在其后他埋头苦干的两年里、乃至之后,我再也没有见过他。他的名义导师是Contou-Carrère,当然;但有机会同这样一个机灵的人聊聊,本会使我高兴。实际上,连答辩的通知都不曾发我,而且我相信从未收到过那份论文的副本——但我记得亲手捧读过一份,是在一个有资格得到它的人手里\fnote{20}{(*)}{这场捉迷藏,愈发离奇,因为毫无疑问:在整个朗格多克,能与Contou-Carrère一道、对Olivier~Leroy的工作多少懂点什么的人,只有我。不必说,Leroy那份科学院报告的底稿,我也同样不曾亲手捧读。也许我在自欺;但我觉得:倘若我不曾被如此苛刻地隔离、以致实际上无从干预,我总归能找到办法,让那篇倒霉的短文发表——必要时经由Cartan或Serre;他们虽不在局中,却肯信任我,只要我担保工作的严肃性。(6月7日)想必很久之后我才知道Leroy已答辩,而且我太忙,无暇自问为何连通知都没人给我。直到Contou-Carrère本人的答辩之后,我才“翻然醒悟”——那场答辩,我被认定是其论文导师(x)。他想方设法,使评审委员会成员中唯独我没有资格得到其论文的正式定本!今天我终于收到一份副本——他以为(他写道)得到它“对我没有兴趣”……(x)更确切地说:有一两年,Contou-Carrère谨慎地同时耍着两位“导师”(谁也说不准……),两人彼此不知有一位“平行”导师。Verdier的导师角色,我是在最后关头才知道的:1983年春,当Verdier分明终究要他的命变得清楚时,Contou-Carrère才最终转而投奔我。}。答辩是在那篇科学院报告“沉没”之前还是之后,我说不好。我说的是那次沉没:在第33节“注记——或新伦理(1)”里,我相当周详但不点名地谈到它。经手这次沉没的两位数学家,是Pierre~Cartier(正是那位:其惊人的直觉之敏捷,在我谈起他年轻的同行的敏捷时重回我心头;而他却如此客气、怀着全部的遗憾,沉没了那份敏捷)与Pierre~Deligne——带着他那句历史性的话:这些数学“不逗他乐”。(它们在他年轻时,可是“逗过乐”的……)我还该加上Contou-Carrère本人:他没有为自己的学生抬一根手指——那会冒触怒权势人物的危险。他想必向Olivier~Leroy暗示:最好忘掉那篇不合时宜的论文。无论如何,清楚的是:Leroy分明对这一情节画了粗粗的一道——即便日后有机会发表的不只是一篇科学院报告、而是其工作的全部,我深信他也不会用\fnote{21}{(*)}{这道“粗线”的一个雄辩征兆:Olivier~Leroy两年前蒙彼利埃一个助教职位出缺时提交的申请材料里,既不提其第三阶段论文的题目,也不提曾是其导师的Contou-Carrère的名字。他也不提任何一项个人的工作,无论什么。显然,他当时还没决定要不要那份职位——于是,尽管其天赋惊人,职位给了另一位候选人:材料扎实、意图毫无疑义的一位。}。这一回,电锯又干得漂亮\fnote{22}{(**)}{一个有趣的巧合:最近我得到回音,Cartier有意把其一场Bourbaki报告献给我(我相信这是头一遭),而且那份报告恰恰专论topos理论——正是这些topos,被同一位Cartier判定不配进入一篇科学院报告。这莫非是最近这几年时尚风向的转变?肯定不是;而且一切都环环相扣:那份报告谈的是topos在逻辑中的使用!我的朋友Cartier这份感人的题献,在我看来,与去年在一个重大场合所致的葬礼悼词同一血脉(见“葬礼悼词——或恭维”一注,n\textsuperscript{o}~104):在那里,“topos”一词被提到(在一串抖得利落的恭维之中),紧接着急忙补上(作为唯一而雄辩的评语)它们“如今用于逻辑”——别处没有,自不待言,至少只要我那些乐善好施恭维的朋友们、凭其手中的权力,还能阻止的话……(Cartier报告的出处:Catégories, logiques et faisceaux, modèles de la théorie des ensembles, Bourbaki研讨班n\textsuperscript{o}~513,1978年2月。)我在某些人(那些定调子的人,如Deligne、Cartier、Quillen……)对topos在几何中这类革新而深刻的观念的那份居高临下的(与抵制的……)态度里,感到一种惊人的狂妄。且假定他们中哪怕一位有那份材料(或那份天真……),能像我经由引入平展topos与晶状topos、从虚无中引出代数簇的一种新的拓扑视野(并由此,在拓扑到来之前,引出代数几何与算术深刻革新的手段)——毫无疑问,他乐于在自己身上培养、在他人身上煽动的这份轻蔑的态度本身,就会解除那视野与革新的力量,只为一份自负效劳。}。

\dmark{}\origpage{412}尽管有这场厄运,1981年初的那几个月里,我到底有幸常见Leroy。那是一场我当时主持的微研讨班,论Teichmüller塔的代数-算术理论(《纲领概要》里略有谈及)。真正意义的听众,只有Contou-Carrère与Leroy。即便对一个超选的巴黎听众(我懂我在说什么),一整间屋子也未必有三四个不掉队的。老实说,我办那场研讨班——当时Contou-Carrère正全力打磨其Schubert闭链的想法——是为了Leroy,想着他或许会对如此辉煌的题目咬钩。分明他“感觉得到”我在做什么;但他(我相信)事先已决定不“咬钩”。奇怪的是他甚至肯费心来——某种东西必定使他着迷,一如我着迷;而他自己对究竟要什么,也不甚了然。当我明白他不会咬钩,我便止损了。面对两位观众(无论多么光彩)继续独白,我无意。我与Leroy的第二场也是最后一场交谈,就在那时。我相信,此后我再没有见过他。

Leroy与我之间,除七年前的那个晚上,没有过真正的数学讨论——这解释了为何对他的工作我几乎一无所知,除那份不合时宜的topos工作。他的厄运想必不曾增进他对像我、甚至Contou-Carrère或“上流”数学界其他人等的信任。我听说他在文学院办研讨班——那里有一群和气的、彼此相处融洽的数学人。他们讲授组合拓扑的想法——一个近十年来全在我射程之内的题目。既然我天性谨慎(可不是嘛!),我没有问他讲些什么,也不知道他是否打算发表。至于境况,他过着一种最不合法的生活(却既非外国人、亦无身份问题):到处带习题课,报酬(嘘……)不知出自哪些小金库,背着司库与审计法院。我想,他是否最终走数学这一行,尚未下定决心;长远看,无论有没有审计法院,这想必不是个舒坦的处境。倘若我那幅富于教益的葬礼图景——他在其中充当第四口编外棺材——能帮他驱散困惑,这一回是在知根知底的情况下,我会高兴。

\subsection{掘墓人——或全体教众}

\notehead{97}{}

(5月24日)我是顶着我身上某种迟疑,才终于决定点名提及我的某些亲近的旧友与同行——在数学界;我得以看到他们充当“掘墓人”(或“电锯”)的角色:从一开始就斩断某些身份寒微或无保障的数学家的尝试——他们试图重提我的某些观念、循其固有的逻辑加以发展,或者(如Yves~Ladegaillerie的个案)仅仅追随一种带着我的影响之印记的进路与做派。如我一说再说,这类牵连他人的迟疑——或仅仅是在未曾\dmark{}\origpage{414}征询其意见的情况下点其名的迟疑\fnote{23}{(*)}{比如,为收进一条注(注n\textsuperscript{o}~19)——其中将点名提及所有跟我做完国家博士论文的学生——我就有过这样的迟疑。我身上这份犹豫,想必来自我的许多学生对与我这个人的名字相连的抗拒;那份抗拒,想必几年前我已在未成形的层面知觉到了。我的旧学生们(带或不带引号)中,当时与我划清界限的意愿已被我清晰知觉的,只有Contou-Carrère(我恰在那时才在他身上发现它)与Deligne(在那里,此事自1968年起已相当分明,尽管我不曾疑心这份意愿将把他带向何方)。在Deligne的个案里,我迟迟不肯以“多少算个学生”的名义点名他的迟疑,格外强烈:我不愿显得在拿这样光彩的学生夸耀,而他自己则不肯让那把他与我、与我的著作连在一起的纽带显露分毫。我的反思,还使我进而明白:这条纽带,在我年轻朋友的生命与著作里所占的分量,远比我曾经猜想的要大。(6月1日)关于我身上这些蓄意的心境,见“三天之后”的3月27日(原书日期如此,与5月24日不合)的“格格不入者”一注(n\textsuperscript{o}~67$'$)。}——在《收获与播种》里并不少见。在每一个个案里,我最终检视了那份迟疑,并懂得它并无根据;其来源不是体贴,而是混乱,甚至怯懦。在我点名陈述他人的行为或态度的所有情形里(依我看),这些行为或态度绝非“保密”性质。它们关涉当事人的职业生活,连同其必然的种种后果——牵连其他同行(包括我自己)的职业生活(并由此,牵连他们的全部生活)。我点名的每一个人,对其行为与态度、及其含义的全部谱系(无论他乐于忽视与否),都与我对我自己的负有同样的责任。倘若其行为的某些后果以这样或那样的形态回到他身上——比如经由我的这个人在此充当媒介的一份公开的“问责”——他毫无根据感到被冒犯。倘若我的语言有时形象而辛辣,我的意图绝非论战,绝非冒犯或羞辱任何人,而宁可是描述事实、描述我感受它们的方式——作为一种敦促,促请每一个人(首先是我点名的那几位)从其自身一侧检视它们,而不是以这样或那样的方式把它们打发掉(一如《收获与播种》的反思之前我常常做的那样)。倘若被如此唤起的人选择感到被冒犯,那是关涉他自己的选择。这份选择可能使我难过——出自我所敬重、甚至怀有感情的人;但它不压垮我。我谈过的那份迟疑——我的事事观感中某种混乱的征兆——一经理解、并由此超越,便无踪无影。

在葬礼的反思里,没有任何时刻,我曾感到有一场针对我的著作、针对那些胆敢从中汲取灵感(而非仅仅借用工具、却讳言锻造并交到他们手里的工匠之名)的人而精心策划的宏大“阴谋”。没有阴谋;但有一个\textbf{共识}——在我所称的“上流”数学界里,它至今在我眼前毫无破绽。这个共识,除至多极罕见的例外,绝非由一份对我这个人、对我的著作的有意识的恶意喂养\fnote{24}{(*)}{我所认为的毫不含糊的恶意行为,我只在Deligne与Verdier的个案里得知。}\dmark{}\origpage{415}。仅在若干例外的个案里,它才表达为对前面各注所谈的四位“共同被埋葬者”中这一位或那一位的、毫不含糊的恶意。但肯定,这样的恶意之所以能在我昔日的某位学生身上滋蔓,之所以能不受阻拦地表达,全赖普遍共识的怂恿。

这份共识,在我的几乎全部旧友或旧学生身上,显现的不是“恶意”的态度,而是(我相信)全然无意识的机制:一种令人错愕的整齐划一、一种毫无破绽的效率,像扫麦秸一样扫走健全意识与数学家的健全本能,让位于纯然自动的\textbf{拒斥态度}\fnote{25}{(**)}{这些拒斥态度,当然从不以这样的面目出现——即便在我朋友Deligne或Verdier那样极端的个案里。在对我有意识的层面,它们几乎不可见;那个层面(如我已有机会说过的)在我昔日的朋友与学生身上,几乎总是(也许甚至总是)同情与敬重的心境(其中某位,还在尽力设法自我开脱)。这样的同情与敬重的心境,不只在有意识的“意见”的表浅层面在场,甚至在实际的好恶、以及对他人实际认识的更深层面上也在场(独立于人们力图把他人关进其中的那些形象)。我们在这里处在一个典型的\textbf{两价性}局面里(我几乎想说:集体性的两价性),粗看之下,什么也“看不见”!(可与“敌父(1)、(2)”(第29、30节)里的反思对读;在《收获与播种》里,我头一回在那里触及这个标记过我生命中许多关系的两价性侧面——不只在数学界。)然而,在具体显现的层面上(葬礼里已被充分检视),这些两价力量的“合力”,在我看来,不再有任何两价性;它分明带着“令人错愕的整齐划一与毫无破绽的效率”,以我正要进一步检视的“自动拒斥态度”的面目呈现。}。这样的自动态度,我疑心,不只由我这个人、由那些其数学“气味”多少让人想起我的人所触发——也对着任何不像已经持有某个“建制”的\textbf{心照不宣的担保}的数学家:无论他已身在其中的行列,还是俨然其中某一位的“门生”(借用Verdier笔下的说法)。我觉得\dmark{}\origpage{416}:在几乎全体数学家身上,最低限度的“数学开放”的心境——它是健全意识与数学健全本能得以启动的前提——\textbf{只对着已经持有这样一份担保的人才启动}。

这类机制,想必几乎是普遍的:不只在数学界,也在社会的所有部门,无一例外。它远远超出一切个案。倘若(依我看来)在我的个案里、以及在建制眼中俨然“我的门生”者的个案里,有某种例外的局面,那是因为:过去,我曾持有“他们自己人”之一的身份,连同对我、对“我的人”的那份惯常的“最低限度的开放”。这份身份,因我1970年的离去而被收回。或者更确切地说:经由我自己的选择——它在我离去后的年月里不止一次清楚地表达,并经由我直至今日的生存方式——我确实不再是一个“他们自己人”。事实上,我自己也不再感到自己是“他们自己人”;我离开了一个曾经共同的世界,绝无回归之念。今天,我的“回归数学”绝不是回归“他们中间”、回归建制,而是回归数学本身;更确切地说,回归一种持续的数学投资,回归对我数学反思的发表活动。

我才开始明白:我的离去,曾被我的旧友们、我的学生们感受为一种“开小差”,甚至一种“冒犯”\fnote{26}{(*)}{这样一种看待与感受事情的方式,在我的朋友Zoghman~Mebkhout身上表达得格外雄辩。正是由于这次开小差,我才要为他与上流数学界的种种挫折负责:唯有他,失去了那份曾环绕着今日乐于把他当拖鞋踢的人的“庇护”与支持。}。想必那是最简单的办法,用以打发我离去的意味、它可能在人们身上唤起的追问:借这样一份弥散的\textbf{受害}之感、借怨恨的自动反应,表达为一场\textbf{报复}之举(在意识的层面,它想必很少被当作报复、甚至当作行为来知觉):既然他与我们切割,我们就与他切割——我们不再给他与“他的人”以保留给“自己人”的“注意的自动机制”的收益——他与他们,将如初来乍到者一般,领受自动拒斥的严苛!

\dmark{}\origpage{417}局面(对我的旧友们与学生们)还更复杂:因为我不但曾是建制的一员,而且他们中任何一个人,要做一个数学家,不可能不随时使用我是作者的观念、想法、工具与结果。我不知道在我们的科学或任何别的科学的历史里,可曾有过如此难堪的悖论的例子!在这一光照下,那些电锯效应(绝不限于我的朋友Deligne)——为斩断任何发展带有我的印记的观念的意图(尽管这样的发展只会加深那份困惑)——如今在我眼前,显出被一种无情的内在逻辑所推动:自某个已然做出的选择——拒斥的选择——起的一份\textbf{必然}。同样,我看到人们四处做出的努力——把那些已进入公共遗产、人们无论愿意与否都离不开的观念、想法、工具与结果的来源,置于彻底的沉默之下——也属同一逻辑。我在Deligne那些非常扎眼的“战役”面前——他装作一项一项地把我的若干主要数学贡献的父权据为己有(或把残渣慷慨地记在某个形影不离的好哥们名下)——曾以为察觉的那份“漠然”——那绝不是漠然,而是一份\textbf{心照不宣的首肯}。Deligne不过是做着建制的集体无意识期待他做的事:\textbf{抹去}那个与所有人切割的人的名字,从而化解那个无法容忍的悖论——\textbf{以一份可容忍的假父权,替换一份真实却不可接受的父权}。

在这一光照下,主祭Deligne,显出的不再是一个按其生命的深层力量与行为塑造一场时尚的人,而分明是\textbf{工具}——凭借其“合法继承人”的角色天造地设的工具:服务于一份毫无破绽的\textbf{集体意志};那意志,执意于一项不可能的任务:把我的名字、我的个人做派,从当代数学中抹去。

我几乎不怀疑:这幅图景,就其实质而言,表达了事情的实在——至少在集体的层面上。肯定,我的“归来”——它出人意料地为一场曾进行得让所有人称心的下葬画上句点,或(即便不画句点)至少不合时宜、不可容忍地扰乱一场仿佛早已排定的仪式的进行——这份归来,将不只使主祭者中的某一位不适与不满,还将使为这场可悲场合齐聚的\origpage{418}全体教众为难!当然,我毫无概念:那鼎鼎大名的集体无意识,将发明怎样的“解数”,来打发这位可敬的亡者不期归来所造成的烂摊子——他从为其量身准备的舒适的棺材里突然爬出(不可容忍的丑闻),装模作样按他自己的方式主持自己的葬礼。但我信得过全体教众:他们会找出办法,打发掉数学大厦里这桩多出来的小小矛盾——它已经万事俱备,只差这个!

我如今,在每个人的意象与态度的层面上,已能相当清楚地知觉集体共识的反映及其总形态,知觉那份抹去、埋葬的集体意志。那是普遍使用的“两幅图景”体系——两幅相互矛盾的图景,人们同时依凭它们行事;我在《收获与播种》里,就我自己的个案,头一回谈到它。(见“功劳与轻蔑”一节,第12节。)我怀疑可有人会直言不讳、明明白白地说:“Grothendieck做的尽是冒牌数学,不提也罢,说正事吧。”照直说,那会太过分明地违背建制的公理——至少眼下如此。按事情预定的演变,二三十年后,这问题反正无从提出:届时连这个名字都无从提起——它早已被所有人遗忘。共同的战术——无论个人还是集体——是沉默:人们不想亡者,至少不作为数学家想;人们不谈他,不提他(除非不得已,借天赐的缩写SGA或EGA;且待这些引证被另一些没有亡者任何踪迹的引证替换……)。

然而,也有——想必例外——的场合,完全的沉默变得行不通。其中一个场合,我料想,曾是我的CNRS任职申请;它想必使不止一人为难\fnote{27}{(*)}{(5月26日)今天我恰好从Zoghman~Mebkhout的电话里得知:我在CNRS全国委员会的同事们为我出了一把力,为我安排了一个两年期的“接待职位”。我不知道他们是否怀着热情做的;总之,委员会里我的朋友中,没有一位把这份出力做到底——打来电话、捎来一句好话,把5月15日就已成定局的好消息告诉我。(9月)我到底经由CNRS一封8月16日的信得到通知:那是一项attaché de recherches职位,为期一年(而非两年)。}。另一场,将是《收获与播种》的预分发\fnotevariant{28}{(*)}{}{说的是由我的大学经手的一份限量印制(150份),在最近的同事与朋友中间分发。},且待其在《数学反思》第一卷里出版(倘若我的出版商不至于胆怯退缩,肯冒得罪整个科学建制的风险)。这些都是亡者本人不可容忍的出格所造成的场合:他不知好歹地走出了派给他的角色。另一个场合(对理解这场尚未被一个不守纪律的亡者搅乱的葬礼,也许更有教益),是去年“隆重”庆祝的IHES二十五周年庆典。作为“IHES的四枚菲尔兹奖章之首”,在这个庄严的场合,要把我全然略过,是困难的——尽管我在IHES存在的四个英雄年月里、为使它取得真实存在所扮演的角色,被略过了。为向我致敬而精心炮制的《葬礼悼词》——载于为此庆典印行的小册子(我已有机会两次指涉它)——依我看是这一类文体的典范:一种优雅而含蓄的方式,使人人满意地化解当代数学里的这份“小小矛盾”……

瞧,我忽然精神焕发——像开始嗅到马厩的马!将近两个星期前,我开始就这一富有教益的情节写一场反思,在一条当即得名“葬礼悼词——或恭维”的注里。经过几番该把它放在哪里的踌躇(它出自为葬礼所写的第一条注的一条迟到的脚注),看来最自然的安放之处,不是“编年”的位置,而是那场应为葬礼压轴的“葬礼仪式”。于是,不经我寻求,我三周来沿着最后三个行列——“研讨会”“学生”,以及刚刚入列的“灵车”(随葬礼的最终部分,即葬礼仪式而来)——所追的那根“线”,接上了头;这场仪式的首要印记,恰恰是那份我5月12日开始检视的《葬礼悼词》杰作\dmark{}\origpage{420},它如今是自然而然紧随本注之后的那条注\fnote{29}{(*)}{(1984年11月)由于一场意外的疾病情节,该注(n\textsuperscript{o}~104)与“本注”被一个新的行列隔开——“亡故者——仍未死去”(n\textsuperscript{o}~98--103)。}。

我终于(又一次?)触到目标了!同时,这场对一篇葬礼悼词的反思的开端,忽然取得一个新的维度。它不再只是一个强力的头脑为一桩执念效劳的精巧发明——在一场官方“重大场合”的贵宾们面前、对着他们的漠然或奉命的关注挥霍——它尤其是一份在所有场合中最微妙的场合、对着一份集体期待所做的完美而体贴周到的回答:那份期待,问的是对我的这个人应当采取的态度。倘若他那一代人中有人真正配得上全体教众毫无保留的感戴,那正是我的朋友Pierre~Deligne——以他那份绝无瑕疵的完美,履行了人们期待于他的角色。
